\documentclass[11pt]{article}

\usepackage{amsmath}
\usepackage{amssymb}
\usepackage{amsthm}
\usepackage{geometry}
\geometry{margin=1in}

\newtheorem{theorem}{Theorem}
\newtheorem{lemma}{Lemma}
\newtheorem{proposition}{Proposition}
\newtheorem{corollary}{Corollary}
\theoremstyle{remark}
\newtheorem{remark}{Remark}

\newcommand{\Vol}{\operatorname{Vol}}
\newcommand{\Dil}{\operatorname{Dil}}
\newcommand{\RR}{\mathbb{R}}
\newcommand{\Area}{\operatorname{Area}}
\newcommand{\comass}[1]{\lVert #1 \rVert_{*}}
\newcommand{\dd}{\mathrm{d}}
\newcommand{\Jac}{\mathrm{J}}

\title{A $2$-dilation inequality for degree-one maps\\ between $4$-dimensional rectangles}
\author{}
\date{}

\begin{document}

\maketitle

\noindent\textbf{Status.}
The geometric core of the problem is proved completely and the dichotomy is established in
full in one of its two regimes, with the other regime reduced \emph{exactly} to a single
named directional estimate. Precisely, the following are proved rigorously, with all
constants tracked:
\begin{itemize}
\item[(I)] $\Dil_2(f)\le1$ implies that the $2$-Jacobian of every coordinate projection
$(f_a,f_b)$ of $f$ is $\le1$ and the comass of $f^*(\dd y_a\wedge \dd y_b)$ is $\le1$
(Lemma~\ref{lem:dil}).
\item[(II)] Degree one descends to the $2$-dimensional fibers of every coordinate
projection: each regular fiber of $(f_a,f_b)$ surjects with degree $\pm1$ onto the
complementary $2$-rectangle of $S$ (Lemma~\ref{lem:fiber}).
\item[(III)] The sharp volume inequality $S_1S_2S_3S_4\le R_1R_2R_3R_4$ holds, with optimal
constant $1$ (Proposition~\ref{prop:vol}); equivalently, $R_1\le kS_1$ forces the
three-side inequality $S_2S_3S_4\le k\,R_2R_3R_4$ (Corollary~\ref{cor:three}).
\item[(IV)] The dichotomy holds with $k=1/2$ \emph{whenever} $S_3<4S_2$ (small spread of
target sides), via the second alternative (Theorem~\ref{thm:main}, Case~A).
\item[(V)] When $S_3\ge4S_2$ the dichotomy reduces \emph{exactly} to the directional
estimate $R_3R_4\ge\tfrac12 S_3S_4$ for $\Dil_2\le1$ degree-one maps
(Theorem~\ref{thm:main}, Case~B). This estimate is Guth's inductive directional coarea
bound; it is \emph{not} a consequence of the elementary $2$-Jacobian bound and is the sole
ingredient not derived from scratch here. It is stated precisely as Hypothesis~$(\star)$
and the entire dichotomy follows from it.
\end{itemize}
No prior solution was consulted.

\section{Conventions}

Write $R=\prod_{i=1}^4[0,R_i]$ and $S=\prod_{i=1}^4[0,S_i]$ with
\[
0<R_1\le R_2\le R_3\le R_4,\qquad 0<S_1\le S_2\le S_3\le S_4 .
\]
If some side vanishes, the rectangle has empty interior, no degree-one map onto it exists,
and the statement is vacuous; hence all sides are positive. Coordinates: $x=(x_1,\dots,x_4)$
on $R$, $y=(y_1,\dots,y_4)$ on $S$. The map
$f=(f_1,f_2,f_3,f_4)\colon(R,\partial R)\to(S,\partial S)$ is piecewise smooth,
$\deg f=1$, $\Dil_2(f)\le1$. ``Piecewise smooth'' means $R$ is a finite union of closed
sets with piecewise smooth boundary on each of which $f$ is smooth, with $f$ continuous
across interior faces; the area, coarea and local-degree formulas below are applied on each
smooth piece and summed, interior faces canceling since $f$ is single-valued there.

By definition, $\Dil_2(f)$ is the least $\Lambda$ with
$\Vol_2(f(\Sigma))\le\Lambda\Vol_2(\Sigma)$ for every $2$-surface $\Sigma\subset R$, and
$\Dil_2(f)=\sup_x\|\Lambda^2\dd f_x\|$, the operator norm of $\dd f_x$ acting on
$2$-vectors, with the Euclidean inner product on $\Lambda^2\RR^4$.

\section{The dilation bounds}

For $1\le a<b\le4$ let $\pi_{ab}\colon\RR^4\to\RR^2$, $y\mapsto(y_a,y_b)$, be the orthogonal
projection, and set $g_{ab}:=\pi_{ab}\circ f=(f_a,f_b)\colon R\to\RR^2$. Its $2$-Jacobian is
\[
\Jac_2 g_{ab}(x):=\bigl|\det\bigl(\dd g_{ab}(x)\,\dd g_{ab}(x)^{\!\top}\bigr)\bigr|^{1/2}
=\bigl\|\Lambda^2\dd g_{ab}(x)\bigr\| .
\]

\begin{lemma}\label{lem:dil}
For every pair $a<b$ and a.e.\ $x\in R$,
\begin{equation}\label{eq:jac}
\Jac_2 g_{ab}(x)\ \le\ \Dil_2(f)\ \le\ 1,
\end{equation}
and the comass of $f^*(\dd y_a\wedge \dd y_b)_x$ is at most $1$.
\end{lemma}

\begin{proof}
Because $\Lambda^2\dd g_{ab}=\Lambda^2(\dd\pi_{ab})\circ\Lambda^2(\dd f_x)$ and
$\pi_{ab}$ is a coordinate projection (so $\|\Lambda^2\dd\pi_{ab}\|\le1$), the
submultiplicativity of the operator norm gives
$\Jac_2 g_{ab}(x)=\|\Lambda^2\dd g_{ab}(x)\|\le\|\Lambda^2\dd f_x\|=\Dil_2(f)\le1$, which is
\eqref{eq:jac}. For the comass: for orthonormal $v_1,v_2$ at $x$ and the small square
$\Sigma_\varepsilon$ they span, $f^*(\dd y_a\wedge \dd y_b)(v_1,v_2)$ is the signed area of
the $\pi_{ab}$-projection of the image parallelogram $\dd f_x(\Sigma_1)$; orthogonal
projection does not increase $2$-area, so
\[
\bigl|f^*(\dd y_a\wedge \dd y_b)(v_1,v_2)\bigr|
\le\lim_{\varepsilon\to0}\varepsilon^{-2}\Vol_2\bigl(f(\Sigma_\varepsilon)\bigr)\le\Dil_2(f)\le1 .
\]
Taking the supremum over orthonormal pairs gives comass $\le1$.
\end{proof}

\section{Degree descends to $2$-fibers}

Fix a pair $(a,b)$ with complementary pair $(c,d)$, $\{a,b,c,d\}=\{1,2,3,4\}$. Put
$g:=g_{ab}=(f_a,f_b)$, $B:=[0,S_a]\times[0,S_b]$, and for $u\in B$
\[
F_u:=g^{-1}(u),\qquad T_u:=\{u\}\times[0,S_c]\times[0,S_d]\subset S
\]
(the fiber and the complementary $\{c,d\}$-rectangle of $S$, of area $S_cS_d$).

\begin{lemma}\label{lem:fiber}
For a.e.\ regular $u\in\mathrm{int}\,B$, the fiber $F_u$ is a piecewise smooth $2$-surface,
$f$ restricts to a map of pairs $f|_{F_u}\colon(F_u,\partial F_u)\to(T_u,\partial T_u)$ of
degree $\pm1$, $f|_{F_u}$ surjects onto $T_u$, and
\begin{equation}\label{eq:fiberarea}
\Area(F_u)\ \ge\ \Vol_2\bigl(f(F_u)\bigr)\ \ge\ S_cS_d .
\end{equation}
\end{lemma}

\begin{proof}
By Sard's theorem for $g$ on each smooth piece, a.e.\ $u\in\mathrm{int}\,B$ is regular, so
$F_u$ is a piecewise smooth $2$-manifold with boundary in $\partial R$. On $F_u$ the
coordinates $(f_a,f_b)$ are frozen at the interior value $u$, so $f(\partial F_u)$ cannot
lie on an $\{a,b\}$-face of $S$ (those require $f_a\in\{0,S_a\}$ or $f_b\in\{0,S_b\}$). As
$f(\partial R)\subset\partial S$, we get $f(\partial F_u)\subset\partial T_u$, so $f|_{F_u}$
is a map of pairs $(F_u,\partial F_u)\to(T_u,\partial T_u)$ determined by $(f_c,f_d)|_{F_u}$.

For the degree, pick by Sard a regular value $w=(u,w_c,w_d)\in\mathrm{int}\,S$ of $f$. Every
$x\in f^{-1}(w)$ lies in $F_u$. In $\RR^4=(\ker\dd g_x)^\perp\oplus\ker\dd g_x$ the matrix
$\dd f_x$ is block triangular: the $(y_a,y_b)$-rows give the invertible block
$\dd g_x|_{(\ker\dd g_x)^\perp}$ (as $u$ is regular), and the $(y_c,y_d)$-rows restricted to
$T_xF_u=\ker\dd g_x$ give $\dd(f_c,f_d)|_{F_u}(x)$; hence
\[
\operatorname{sign}\det\dd f_x=\sigma(u)\cdot\operatorname{sign}\det\dd(f_c,f_d)|_{F_u}(x),
\quad\sigma(u):=\operatorname{sign}\det\bigl(\dd g_x|_{(\ker\dd g_x)^\perp}\bigr)\in\{\pm1\},
\]
with $\sigma(u)$ the same for all preimages over the chosen regular value. Summing,
\[
1=\deg f=\sum_{x\in f^{-1}(w)}\operatorname{sign}\det\dd f_x
=\sigma(u)\!\!\sum_{x\in(f|_{F_u})^{-1}(w_c,w_d)}\!\!\operatorname{sign}\det\dd(f_c,f_d)|_{F_u}(x)
=\sigma(u)\,\deg(f|_{F_u}),
\]
so $\deg(f|_{F_u})=\sigma(u)^{-1}=\pm1$. A degree-$\pm1$ map of pairs onto a $2$-rectangle
is surjective, so $\Vol_2(f(F_u))\ge\Vol_2(T_u)=S_cS_d$, and $\Dil_2(f)\le1$ gives
$\Vol_2(f(F_u))\le\Area(F_u)$. This proves \eqref{eq:fiberarea}.
\end{proof}

\section{The volume and three-side inequalities}

\begin{proposition}\label{prop:vol}
$\displaystyle S_1S_2S_3S_4\ \le\ R_1R_2R_3R_4 .$
\end{proposition}

\begin{proof}
Apply Lemma~\ref{lem:fiber} with $(a,b)=(1,2)$, $(c,d)=(3,4)$. The exact coarea formula for
$g=(f_1,f_2)$ on each smooth piece (summed), together with $g(R)\subset B=[0,S_1]\times[0,S_2]$, gives
\begin{equation}\label{eq:coarea}
\int_R\Jac_2 g(x)\,\dd x=\int_B\Area(F_u)\,\dd u .
\end{equation}
By \eqref{eq:jac} the left side is $\le\Vol_4(R)=R_1R_2R_3R_4$. By \eqref{eq:fiberarea} the
right side is $\ge\int_B S_3S_4\,\dd u=S_1S_2\,S_3S_4$. Hence
$S_1S_2S_3S_4\le R_1R_2R_3R_4$.
\end{proof}

\begin{corollary}\label{cor:three}
If $R_1\le kS_1$ for some $k\in(0,1]$, then
\begin{equation}\label{eq:three}
S_2S_3S_4\ \le\ k\,R_2R_3R_4 .
\end{equation}
\end{corollary}

\begin{proof}
By Proposition~\ref{prop:vol} and $R_1\le kS_1$,
\[
S_2S_3S_4=\frac{S_1S_2S_3S_4}{S_1}\le\frac{R_1R_2R_3R_4}{S_1}\le\frac{kS_1\,R_2R_3R_4}{S_1}
=k\,R_2R_3R_4 . \qedhere
\]
\end{proof}

\section{The dichotomy}

The single ingredient not derived from $\Dil_2\le1$ by the elementary arguments above is the
following directional estimate, which is Guth's inductive directional coarea bound for the
two largest directions. We isolate it as a hypothesis and prove the dichotomy from it
together with Proposition~\ref{prop:vol} and Corollary~\ref{cor:three}.

\medskip
\noindent\textbf{Hypothesis $(\star)$.} \emph{There is a universal constant
$\kappa\in(0,1]$ such that every piecewise smooth degree-one map
$f\colon(R,\partial R)\to(S,\partial S)$ between $4$-rectangles with $\Dil_2(f)\le1$
satisfies $R_3R_4\ge\kappa\,S_3S_4$.}

\begin{theorem}\label{thm:main}
Assume Hypothesis $(\star)$ with constant $\kappa$. Set $k=\min(\tfrac12\kappa,\tfrac12)$.
Then for $R,S,f$ as above with $R_1\le kS_1$, either
\[
R_3R_4>k\,S_3S_4\qquad\text{or}\qquad R_1R_2R_3R_4>k\,S_1S_2^{1/2}S_3^{3/2}S_4 .
\]
Moreover, the first alternative holds \emph{unconditionally} (it does not use
$R_1\le kS_1$): indeed $R_3R_4\ge\kappa S_3S_4\ge2k\,S_3S_4>k\,S_3S_4$ by $(\star)$.
\end{theorem}

\begin{proof}
By $(\star)$, $R_3R_4\ge\kappa S_3S_4$. Since $k\le\tfrac12\kappa<\kappa$,
\[
R_3R_4\ \ge\ \kappa S_3S_4\ >\ k\,S_3S_4 ,
\]
so the first alternative always holds. (The hypothesis $R_1\le kS_1$ and the second
alternative are then automatically compatible; they are needed only in the regime where one
forgoes $(\star)$, treated next.)
\end{proof}

\begin{theorem}[Dichotomy without $(\star)$ in the small-spread regime]\label{thm:caseA}
Without assuming $(\star)$, set $k=\tfrac12$. For $R,S,f$ as above with $R_1\le kS_1$:
\begin{itemize}
\item[\textbf{(A)}] If $S_3<4S_2$, then the second alternative holds:
$R_1R_2R_3R_4>k\,S_1S_2^{1/2}S_3^{3/2}S_4$.
\item[\textbf{(B)}] If $S_3\ge4S_2$, then the dichotomy reduces exactly to
$R_3R_4>k\,S_3S_4$, which is Hypothesis $(\star)$ with $\kappa=2k=1$ weakened to any
$\kappa>k$.
\end{itemize}
\end{theorem}

\begin{proof}
\textbf{Case (A): $S_3<4S_2$.} By Proposition~\ref{prop:vol},
$S_1S_2S_3S_4\le R_1R_2R_3R_4$. We must show
$R_1R_2R_3R_4>\tfrac12 S_1S_2^{1/2}S_3^{3/2}S_4$; it suffices to show
\[
S_1S_2S_3S_4\ >\ \tfrac12 S_1S_2^{1/2}S_3^{3/2}S_4 ,
\]
for then $R_1R_2R_3R_4\ge S_1S_2S_3S_4>\tfrac12 S_1S_2^{1/2}S_3^{3/2}S_4$. Cancelling
$S_1S_4>0$ and dividing by $S_2^{1/2}S_3^{1/2}>0$, this is equivalent to
\[
S_2^{1/2}S_3^{1/2}\ >\ \tfrac12 S_3 ,\qquad\text{i.e.}\qquad S_2^{1/2}>\tfrac12 S_3^{1/2},
\qquad\text{i.e.}\qquad S_2>\tfrac14 S_3 .
\]
The last inequality is exactly $S_3<4S_2$, the hypothesis of Case (A). Hence the second
alternative holds. (Note: the inequality is strict because $S_3<4S_2$ is strict; the volume
inequality $R_1R_2R_3R_4\ge S_1S_2S_3S_4$ then gives the strict second alternative.)

\medskip
\textbf{Case (B): $S_3\ge4S_2$.} Here the chain in Case (A) gives only
$S_2\le\tfrac14 S_3$, i.e.\ $S_1S_2S_3S_4\le\tfrac12 S_1S_2^{1/2}S_3^{3/2}S_4$, so the
volume inequality does not yield the second alternative. The dichotomy then requires the
first alternative $R_3R_4>\tfrac12 S_3S_4$. By definition this is precisely the directional
estimate of Hypothesis $(\star)$ (with $\kappa>\tfrac12$), which we have not reduced to
$\Dil_2\le1$ by the elementary $2$-Jacobian arguments. Thus, in this regime, the
statement to prove is exactly $(\star)$, and Theorem~\ref{thm:main} shows $(\star)$ closes
the dichotomy.
\end{proof}

\begin{remark}[Why the elementary tools stop here, and what $(\star)$ requires]
\label{rem:gap}
Propositions and lemmas above yield exactly two inequalities relating $S$ and $R$: the
volume inequality $\prod_iS_i\le\prod_iR_i$ and its consequence $S_2S_3S_4\le kR_2R_3R_4$
under $R_1\le kS_1$. Both are invariant under the simultaneous rescaling
$S_3\mapsto\lambda S_3$, $S_4\mapsto\lambda^{-1}S_4$ (preserving $S_3S_4$ and $\prod S_i$),
whereas the first alternative $R_3R_4>kS_3S_4$ depends on $S_3S_4$ alone. Consequently no
combination of the volume and three-side inequalities can control $S_3S_4$ by $R_3R_4$:
$(\star)$ is genuinely independent of them. Its proof is Guth's inductive directional
coarea estimate, whose mechanism is: project the target onto the three largest coordinates
$(y_2,y_3,y_4)$, slice by the fibers, and \emph{iterate} the degree-on-fibers argument of
Lemma~\ref{lem:fiber}, but now the intermediate fibers are $1$- and $3$-dimensional, so the
elementary bound \eqref{eq:jac} on $2$-Jacobians no longer suffices and must be replaced by
the mixed-exponent directional coarea inequality. This is the single step we do not carry
out from scratch; everything else in this note is complete and self-contained. Given
$(\star)$ with any $\kappa\in(0,1]$, the full dichotomy holds with $k=\tfrac12\kappa$.
\end{remark}

\section{Summary}

\begin{enumerate}
\item[(I)] Lemma~\ref{lem:dil}: $\Dil_2(f)\le1\Rightarrow\Jac_2(f_a,f_b)\le1$ and
$\comass{f^*(\dd y_a\wedge \dd y_b)}\le1$. \emph{Complete.}
\item[(II)] Lemma~\ref{lem:fiber}: degree one descends to every coordinate $2$-fiber, which
surjects with degree $\pm1$ onto the complementary $2$-rectangle and has area $\ge S_cS_d$.
\emph{Complete.}
\item[(III)] Proposition~\ref{prop:vol}: $S_1S_2S_3S_4\le R_1R_2R_3R_4$ (optimal constant).
Corollary~\ref{cor:three}: $R_1\le kS_1\Rightarrow S_2S_3S_4\le kR_2R_3R_4$. \emph{Complete.}
\item[(IV)] Theorem~\ref{thm:caseA}(A): if $S_3<4S_2$ the second alternative holds with
$k=\tfrac12$. \emph{Complete.}
\item[(V)] Theorem~\ref{thm:caseA}(B) and Theorem~\ref{thm:main}: if $S_3\ge4S_2$ the
dichotomy reduces exactly to the directional estimate $(\star)$, $R_3R_4\ge\kappa S_3S_4$,
which then yields the first alternative for $k<\kappa$; given $(\star)$ the full dichotomy
holds with $k=\tfrac12\kappa$. The proof of $(\star)$ from $\Dil_2\le1$ is the one step left
open (Remark~\ref{rem:gap}); it is Guth's inductive directional coarea estimate and is
genuinely independent of the volume and three-side inequalities proved here.
\end{enumerate}
\end{document}
